\section{notable points}

In this section, we will introduce terminology that is widely used in the study
of functions. We will attempt to provide definitions even for some terms on
which there may not be a unanimous consensus. Let us remember not to use these
definitions too formally: it is always better to ensure that our interlocutor
understands us, as often some terms might be used improperly or with slightly
different meanings.

In case of doubt, we can always avoid using this terminology by referring to the
underlying concepts.

\begin{definition}[notable points]
Let $f\colon A \subset \RR \to \RR$ be a function. If $f$ is differentiable at a
point $x_0\in A$ and $f'(x_0) = 0$, we will say that $x_0$ is a \emph{critical
point}%
\mymargin{critical point}%
\index{point!critical}
or
\emph{stationary point}
\index{point!stationary}
of $f$.

If $x_0\in A$ and there exists a neighborhood $U$ of $x_0$ such that $x_0$ is a
point of minimum (respectively maximum) for $f$ restricted to $U$, we will say
that $x_0$ is a point of \emph{relative minimum}%
\mymargin{relative minimum}%
\index{minimum!relative} or \emph{local minimum}
(respectively \emph{relative maximum} or \emph{local maximum}).
In contrast, the points of maximum and minimum over the entire domain $A$ are
also called the \emph{absolute} maximum/minimum of $f$.

We will say that $x_0\in A$ is an \emph{inflection point}%
\mymargin{inflection point}%
\index{point!inflection} for $f$ if
$f$ is differentiable in a neighborhood of $x_0$ and $x_0$ is a point of maximum
or minimum for $f'$. At the point $x_0$, the tangent line has the equation
$y=r(x) = f'(x_0) (x-x_0) + f(x_0)$. If $x_0$ is a minimum for $f'$, it follows
that $f(x)-r(x)$ is increasing, so $f(x)\ge r(x)$ for $x\ge x_0$ and $f(x)\le
r(x)$ for $x\le x_0$ (the graph of the function crosses the tangent line from
below to above), whereas if $x_0$ is a maximum for $f'$, it follows that
$f(x)\le r(x)$ for $x\ge x_0$ and $f(x) \ge r(x)$ for $x\le x_0$ (the graph of
the function crosses the tangent line from above to below).
If the function $f$ is not differentiable at $x_0$ but the limit of the
difference quotient exists and is infinite, we will say that $x_0$ is a
\emph{vertical inflection point}%
\mymargin{vertical inflection point}%
\index{inflection!vertical}. At this point, the tangent line is vertical, and the graph of the function crosses this line.

Let $x_0\in A$ be a point where the function $f$ is continuous and the right and
left limits of the difference quotient exist (which are called the \emph{right
derivative} and \emph{left derivative})
\[
  m^{\pm} = \lim_{h\to 0^\pm}\frac{f(x+h) - f(x)}{h}.
\]
If $m^+ \neq m^-$, clearly $f$ is not differentiable at $x_0$.
If both $m^+$ and $m^-$ are finite, we will say that $x_0$ is a \emph{cusp
point}%
\mymargin{cusp point}%
\index{point!cusp} because the two tangent half-lines at $x_0$ (from the right and from the left) form a non-flat angle.
If $m^+=-m^+=+\infty$ or if $m^+=-m^-=-\infty$, we will say that the point $x_0$
is a \emph{cusp point}%
\mymargin{cusp point}%
\index{point!cusp} (there is a vertical tangent half-line).
\end{definition}

\begin{definition}[asymptotes]
We will say that the line $y=mx+q$ is an asymptote for the graph of $f$ as $x\to
+\infty$ if it holds that
\[
  \lim_{x\to +\infty} f(x) - (mx+q) = 0.
\]
If $m=0$, we will say that the graph of $f$ has a \emph{horizontal asymptote}
$y=q$
\mymargin{horizontal and oblique asymptote}%
\index{asymptote horizontal and oblique}%
\index{asymptote!horizontal}%
\index{asymptote!oblique}%
otherwise, we will say that $y=mx+q$ is an \emph{oblique asymptote}.
The same can be said for $x\to -\infty$.

If $x_0\in \RR$ and it holds that 
\[
  \lim_{x\to x_0} \abs{f(x)} = +\infty
\]
we will say that the line $x=x_0$ is a \emph{vertical asymptote}%
\mymargin{vertical asymptote}%
\index{asymptote!vertical} for the graph of the function $f$.
\end{definition}

\begin{definition}[points of discontinuity]
If for $x_0\in \RR$ it holds that 
\[
  \lim_{x\to x_0^-} f(x) = \ell_1, 
  \qquad 
  \lim_{x\to x_0^+} f(x) = \ell_2
\]
and if $\ell_1\neq \ell_2$, we will say that at the point $x_0$ the function $f$
has a \emph{jump discontinuity}%
\mymargin{jump discontinuity}%
\index{discontinuity!jump}.
If $\ell_1=\ell_2$ and if $f$ is not defined at the point $x_0$ or if
$f(x_0)\neq \ell_1$, we will say that at the point $x_0$ the function $f$ has a
\emph{removable discontinuity}%
\mymargin{removable discontinuity}%
\index{discontinuity!removable}.
\end{definition}

It should be noted that, despite the terminology used, a continuous function can
have a jump discontinuity (for example: $f(x)=\frac{x}{\abs{x}}$) and can also
have a removable discontinuity (for example: $f(x) = \frac{x}{x}$).